\chapter{Moles and Nevi}

\section*{Definition}
A melanocytic nevus (mole) is a benign proliferation of melanocytes at the dermoepidermal junction (junctional), within the dermis (intradermal), or at both sites (compound). Most adults have 10--40 nevi; counts \ensuremath{>}50 increase melanoma risk.

\section*{What Happens in Your Body}
Nevi arise from neural crest-derived melanocytes that migrate to the skin during embryogenesis. Activating BRAF V600E mutations (70--80\% of acquired nevi) drive initial melanocyte proliferation, after which oncogene-induced senescence (p16INK4a upregulation) halts growth, producing benign lesions that remain stable for years. Sun exposure, particularly during childhood, accelerates nevus development.

\section*{Causes}
\begin{itemize}
  \item \textbf{Genetic:} Familial atypical mole and melanoma syndrome (FAMMM, CDKN2A mutations), BRAF and NRAS mutations
  \item \textbf{Environmental:} UV radiation (especially intermittent sun exposure and blistering sunburns in childhood), immunosuppression
  \item \textbf{Hormonal:} Increased number and darkening during pregnancy, puberty, and oral contraceptive use
  \item \textbf{Congenital:} Present at birth or within first 6 months (congenital melanocytic nevi); giant congenital nevi (>20 cm) carry \ensuremath{>}5\% lifetime melanoma risk
\end{itemize}

\section*{When to See a Doctor}
See your doctor if:
\begin{itemize}
  \item ABCDE criteria: Asymmetry, irregular Border, color variegation, Diameter \ensuremath{>}6 mm, Evolution/change over time
  \item Ugly duckling sign: A mole that looks distinctly different from all others
  \item Itching, bleeding, crusting, or ulceration of a pre-existing nevus
  \item New mole after age 30--40 years (warrants careful evaluation)
\end{itemize}

\section*{Related Symptoms}
\begin{itemize}
  \item Melanoma, atypical (dysplastic) nevus, Spitz nevus, blue nevus, halo nevus, lentigo
\end{itemize}
\textbf{Important:} Total body photography and serial dermoscopic imaging every 3--12 months is standard surveillance for people with atypical mole syndrome or strong family history of melanoma.

\section*{Self-Care / Home Management}
\begin{itemize}
  \item Monthly self-skin examination using the ABCDE mnemonic; use a mirror or partner for back and scalp
  \item Daily broad-spectrum sunscreen (SPF \ensuremath{\geq}50) and sun-protective behavior
  \item Track photographs of suspicious or changing moles for comparison over time
\end{itemize}

\section*{How Your Doctor Will Evaluate You}
\begin{itemize}
  \item Dermoscopic examination: Pattern analysis (reticular, globular, starburst, homogeneous, blue-white veil)
  \item Sequential digital dermoscopy imaging for short-term monitoring (3--6 months) of suspicious lesions
  \item Excisional biopsy with 1--2 mm margins for histopathologic diagnosis of any lesion concerning for melanoma
\end{itemize}

\section*{Treatment Options}
\begin{itemize}
  \item Benign nevi: No treatment necessary; cosmetic excision for patient preference or irritation
  \item Atypical nevi: Excision with clear margins and histologic evaluation; annual skin surveillance
  \item Congenital nevi: Consider prophylactic excision for giant nevi; small/medium congenital nevi are observed unless changes occur
  \item Shave removal for raised, dome-shaped intradermal nevi; punch or excision for flat or atypical lesions
\end{itemize}

\section*{References}
\begin{thebibliography}{99}
\bibitem{ref1} Bataille V, Bishop JA, Sasieni P, et al. "Risk of cutaneous melanoma in relation to the numbers, types and sites of naevi." \textit{International Journal of Cancer}, 1996;65(1):16--21.
\bibitem{ref2} Pollock PM, Harper UL, Hansen KS, et al. "High frequency of BRAF mutations in nevi." \textit{Nature Genetics}, 2003;33(1):19--20.
\bibitem{ref3} Tsao H, Bevona C, Goggins W, et al. "The transformation rate of moles (melanocytic nevi) into cutaneous melanoma." \textit{Archives of Dermatology}, 2003;139(3):282--288.
\bibitem{ref4} Marghoob AA, Kopf AW, Rigel DS, et al. "Clinical and dermoscopic features of common and atypical nevi." \textit{Dermatologic Clinics}, 2005;23(3):407--421.
\bibitem{ref5} Garbe C, Burg G, Falanga V, et al. "Diagnosis and treatment of melanoma." \textit{Lancet}, 2016;387(10030):1209--1220.
\bibitem{ref6} Schwartz RJ, Kleinerman RA, Linos E, et al. "Diagnosis and management of melanocytic nevi." UpToDate, 2023.

\end{thebibliography}
